\documentclass[11pt]{article}
\usepackage{mathunal-base}
\mathunalfuente{serif}

\renewcommand{\examtitulo}{Segundo Parcial de Matemáticas III}
\renewcommand{\examfecha}{Semestre 02 de 2007 --- Noviembre 3 de 2007}

\begin{document}

\examheader

\vspace{4pt}
\noindent Nombre: \dotfill\ Carné: \dotfill

\vspace{4pt}
\noindent Profesor: \dotfill

\begin{center}\textbf{Duración: 2 horas}\end{center}

\noindent\textbf{Observaciones:} - Justifique cuidadosamente sus respuestas.

\hspace{2.3cm} - Cualquier tentativa de fraude es causal de anulación del examen.

\vspace{8pt}
\begin{enumerate}
\item{} (20\%) Sea $f:\mathbb{R}^2\longrightarrow\mathbb{R}$ la función definida por $f(x,y)=2x^3-3x^2+y^2+2y$. Halle los puntos críticos de $f$ y determine cuáles de ellos son puntos de mínimo local, cuáles son de máximo local y cuáles son puntos de silla.

\vspace{4cm}

\item{} (25\%) Sea $E$ el subconjunto del primer octante de $\mathbb{R}^3$ limitado por los planos $y=1-x$, $z=0$, $y=0$ y $x=0$ y por la superficie $z=1-y^2$. Escriba la integral triple correspondiente al volumen de $E$ de tres formas distintas usando coordenadas cartesianas y diferentes órdenes de integración.

\noindent \textbf{Nota}: no es necesario que calcule la integral, sólo déjela planteada.

\vspace{5cm}

\item{} (30\%) Sea $\Omega\subset\mathbb{R}^3$ el conjunto definido por
\[
\Omega := \{(x,y,z)\ |\ x^2+y^2+z^2\leq 2;\ z\geq 1\}.
\]
\noindent Suponga que en cada punto de $\Omega$ la densidad de masa es $\sigma(x,y,z)=x^2+y^2+z^2$.

\vspace{4pt}
\noindent a) \textbf{Escriba} la masa de $\Omega$ como una integral \textbf{triple} en las coordenadas \textbf{cilíndricas} $(r,\theta,z)$ especificando los extremos de integración para cada variable. Justifique claramente los extremos de integración correspondientes a la variable cilíndrica $r$.

\noindent \textbf{Nota}: no es necesario que calcule la integral, sólo déjela planteada.

\vspace{4pt}
\noindent b) \textbf{Escriba} la masa de $\Omega$ como una integral \textbf{triple} en las coordenadas \textbf{esféricas} $(\rho,\theta,\phi)$ especificando los extremos de integración para cada variable. Justifique claramente los extremos de integración correspondientes a la variable esférica $\phi$.

\noindent \textbf{Nota}: no es necesario que calcule la integral, sólo déjela planteada.

\vspace{6cm}

\item{} (25\%) Calcule
\[
\iint_D e^{\left(\frac{y}{x}\right)} \frac{(xy)^2}{8-\left(\frac{y}{x}\right)^3}\left(\frac{y}{x}\right)\,dA,
\]
\noindent donde $D$ es la región del plano limitada por las rectas $\frac{y}{x}=\frac{1}{2}$, $\frac{y}{x}=1$ y $x=1$ y por la hipérbola $xy=2$.

\noindent \textbf{Nota}: para simplificar los cálculos se puede usar el hecho de que
\[
\frac{\partial(x,y)}{\partial(u,v)} = \left(\frac{\partial(u,v)}{\partial(x,y)}\right)^{-1}.
\]

\vspace{6cm}
\end{enumerate}

\examfooter
\end{document}
